% ============================================================
%  04_pretraitement.tex — Section 4 : Pre-traitement
% ============================================================
\section{Pré-traitement et réduction du domaine}

Le pré-traitement exploite la structure du graphe dual \textit{avant} la résolution pour
réduire l'espace de recherche. L'idée centrale est que la table de voisinage complète (1203
entrées) est un catalogue \textit{générique}. Une fois le graphe dual du benzénoïde d'entrée
connu, il est possible de ne conserver que le \textbf{sous-ensemble pertinent} de la table pour
chaque hexagone, en fonction de son motif $\sigma(v)$ et de son degré dans le dual. Ce filtrage
réduit considérablement le nombre de combinaisons que le solveur doit explorer.

\subsection{Calcul des domaines et hexagones gelés}

Chaque variable $x_v$ représente la taille du cycle à la position $v$. Le domaine est réduit
selon la nature de l'hexagone. Trois règles indépendantes sont appliquées :

\begin{enumerate}
  \item \textbf{Règle deg=6} (toujours active) : si un hexagone est entouré de 6 voisins,
    il ne peut pas changer de taille. Son domaine est fixé à $\{6\}$.
  \item \textbf{Règle $b(v) \geq 2$} (désactivable) : si les arêtes partagées découpent le
    bord de l'hexagone en deux blocs ou plus d'arêtes libres, la substitution est impossible
    car elle déformerait un seul bloc sans affecter les autres. Le domaine est fixé à $\{6\}$.
  \item \textbf{Filtrage topologique C2bis} (toujours active) : pour les hexagones non
    gelés, on n'ajoute $5$ au domaine que si l'hexagone admet un \emph{sommet intérieur
    libre} (deux arêtes adjacentes libres en succession dans $\sigma(v)$), et on n'ajoute
    $7$ que si l'hexagone a au moins un côté libre. Sans ces conditions, la reconstruction
    3D est topologiquement impossible.
\end{enumerate}

\dateblock{mai 2026 (fin de mois)}

La règle C2bis a été ajoutée après avoir observé sur $h_8$ et $h_9$ une grande proportion
de solutions CSP-valides que la reconstruction rejetait avec \texttt{ValueError}
(\og solutions \texttt{geom\_infeasible} \fg{}) --- jusqu'à 75\,\% sur certaines molécules
de $h_9$ en mode \texttt{-{}-no-freeze -{}-no-table}. Filtrer en amont au niveau du domaine
divise par 3 à 5 le temps de reconstruction sur ces configurations.

\begin{lstlisting}
def has_interior_free_vertex(pattern):
    """Vrai ssi le pattern contient au moins 2 zeros consecutifs (cyclique).
    Condition pour contracter l'hexagone en pentagone."""
    n = len(pattern)
    return any(pattern[(i - 1) % n] == 0 and pattern[i] == 0
               for i in range(n))


def has_free_side(pattern):
    """Vrai ssi le pattern contient au moins un zero (cote libre).
    Condition pour inserer un sommet (passage a l'heptagone)."""
    return any(p == 0 for p in pattern)


def compute_domains(graph, freeze_b2=True):
    domains = {}
    for v in range(graph.h):
        pattern = graph.patterns[v]
        if graph.is_fully_surrounded(v):              # deg = 6
            domains[v] = {6}
        elif freeze_b2 and graph.has_separated_free_edges(v):  # b(v) >= 2
            domains[v] = {6}
        else:
            d = {6}
            if has_interior_free_vertex(pattern):     # C2bis: 5 admissible
                d.add(5)
            if has_free_side(pattern):                # C2bis: 7 admissible
                d.add(7)
            domains[v] = d
    return domains
\end{lstlisting}

\begin{remarque}
  Les deux helpers \texttt{has\_interior\_free\_vertex} et \texttt{has\_free\_side}
  appliquent exactement la même condition que la reconstruction (\texttt{reconstruction/
  topology.py::\_get\_interior\_free\_vertices}), garantissant la cohérence entre le
  filtrage CSP et la phase de reconstruction. Si la reconstruction évolue, ces helpers
  doivent être maintenus en parallèle.
\end{remarque}

\begin{remarque}[Cas où C2bis débloque des solutions]
  Lorsque \texttt{--no-freeze} désactive la règle $b(v) \geq 2$, C2bis reste active.
  Un pattern $(0,0,1,0,1,1)$ (deux blocs, mais l'un contient deux zéros adjacents)
  obtient le domaine $\{5, 6, 7\}$ : C2bis autorise légitimement le pentagone là où
  l'ancienne version se contentait de $\{5, 6, 7\}$ aveuglément, et là où un
  $(1,1,1,1,1,0)$ se voit désormais réduit à $\{6, 7\}$.
\end{remarque}

\begin{lstlisting}[style=benzai]
python main.py data/5.graph              # les trois regles actives
python main.py data/5.graph --no-freeze  # deg=6 + C2bis actives, b>=2 desactive
\end{lstlisting}

\subsection{Filtrage de la table de voisinage}

\dateblock{mai 2026 (mi-mois)}

\textbf{Évolution :} la table a été regénérée avec le protocole MD (cf. section sur
\texttt{non\_benzenoid\_generator/main.py optimize}). Tailles avant/après :

\begin{center}
\begin{tabular}{lrrr}
  \toprule
                & xTB \texttt{--opt} (avr. 2026) & MD + opt (mai 2026) & Variation \\
  \midrule
  $\mathcal{T}(5)$ &  92 &  63 & --32\% \\
  $\mathcal{T}(6)$ & 418 & 280 & --33\% \\
  $\mathcal{T}(7)$ & 699 & 367 & --48\% \\
  \midrule
  \textbf{Total}  & \textbf{1209} & \textbf{710} & \textbf{--41\%} \\
  \bottomrule
\end{tabular}
\end{center}

La baisse de 41\% confirme l'hypothèse : une part non négligeable de l'ancienne table
était des \textbf{faux positifs planaires}, c'est-à-dire des configurations qui semblaient
planes parce que xTB \texttt{--opt} restait piégé dans le minimum 2D plat parasite. Le
protocole MD casse ces minima et ramène les structures à leur géométrie réelle (souvent
courbée). Les entrées rejetées sont donc des cas que le solveur CSP \emph{n'aurait jamais
dû} accepter comme bases de substitution admissibles.

\begin{remarque}
  La table $\mathcal{T}(6)$ inclut toujours les 6 entrées purement benzenoïdes (hexagone
  central avec uniquement des voisins hexagonaux). Elles permettraient de retrouver le
  benzénoïde d'origine comme solution CSP, mais celui-ci est désormais \textbf{exclu par
  defaut} des solutions énumérées (cf. drapeau \texttt{-{}-count-hexagon} en section~5).
\end{remarque}

Pour chaque hexagone libre $v$ de degré $\geq 2$, on filtre la table pour ne garder que les
entrées \textbf{compatibles} avec le patron $\sigma(v)$ :

\begin{lstlisting}
def filter_table_for_vertex(graph, v, full_table):
    pattern = graph.patterns[v]
    deg = graph.degree(v)
    shared_positions = [i for i, p in enumerate(pattern) if p == 1]

    compatible = []
    for n in (5, 6, 7):
        for seq in full_table[n]:
            seq_shared = [i for i, s in enumerate(seq) if s != 0]
            # Le nombre de voisins doit correspondre
            if len(seq_shared) != deg:
                continue
            # Les positions non-nulles doivent former un bloc consecutif
            if not _is_consecutive_block(seq_shared, n):
                continue
            # Construire le tuple pour la contrainte extensionnelle
            neighbor_sizes = [seq[i] for i in seq_shared]
            compatible.append((n,) + tuple(neighbor_sizes))

    return compatible
\end{lstlisting}

Le filtrage vérifie deux conditions :
\begin{enumerate}
  \item Le nombre de voisins non nuls dans la séquence correspond au degré de $v$ dans le dual.
  \item Les positions non nulles forment un \textbf{bloc consécutif} (cyclique) — condition
    nécessaire puisque $b(v) = 1$ pour les hexagones libres.
\end{enumerate}

\begin{remarque}
  Les hexagones de degré 1 (un seul voisin) ne nécessitent pas de contrainte de table :
  un cycle avec un unique voisin est toujours géométriquement admissible, quelle que soit
  sa taille.
\end{remarque}

\subsection{Calcul des automorphismes}

Les \textbf{automorphismes} du graphe dual sont utilisés pour la brisure de symétrie
(contrainte C4). Un automorphisme est une permutation des sommets qui préserve la structure
d'adjacence.

\begin{lstlisting}
def compute_automorphisms(graph):
    G = graph.dual
    matcher = nx.algorithms.isomorphism.GraphMatcher(G, G)
    # Collecter tous les automorphismes (sauf l'identite)
    all_auts = [iso for iso in matcher.isomorphisms_iter()
                if any(iso[v] != v for v in iso)]
    # Extraire un ensemble minimal de generateurs
    return _extract_generators(all_auts, list(G.nodes()))
\end{lstlisting}

L'extraction des générateurs utilise une approche gloutonne : on ajoute un automorphisme au
jeu de générateurs s'il produit de nouvelles permutations par composition avec les générateurs
existants. On s'arrête quand la fermeture du sous-groupe couvre tous les automorphismes.
